\documentclass[11pt]{article}
\usepackage[top=.64in,bottom=.69in,left=.8in,right=.8in]{geometry}
\usepackage{amsmath,amssymb,amsthm,booktabs,microtype,xcolor}
\usepackage[colorlinks=true,urlcolor=blue!45!black]{hyperref}
\newtheorem{theorem}{Theorem}
\linespread{.95}
\newcommand{\F}{\mathbb F}
\newcommand{\PP}{\mathbb P}
\begin{document}
\begin{center}
{\LARGE Three- and four-step graph towers}\\[3pt]
{\large Distributions and extremizer orbits}\\[6pt]
{\small 27 November 2026}
\end{center}
\begin{abstract}
We give the full three-step valuation distribution, including a
classification of the three extremizer orbits when $p\equiv1\bmod6$.
For four steps with identically zero quadratic moment we compute the
distribution from a finite-field curve point count and classify the
extremizers when $p\equiv1\bmod8$. The four-step maximum itself is
already known from the cited ProveIt report.
\end{abstract}

\section{Setup and three steps}
Fix a prime $p>2m$ and a spanning list of $m$ step residues in
$\F_p^2$. Form the $2\times m$ matrix of columns, with row plane
$W\subset\F_p^m$. Changing the basis in $\F_p^2$ gives precisely
the other matrices with row plane $W$. We count these row planes
before quotienting by sign changes and permutations.
For a direction $L\in\PP(W)$ let $r(L)$ be the least $j\ge1$
for which the sum of the $2j$th powers of a representative's
coordinates is nonzero. If $A=2\sum_L r(L)$, the
\href{https://github.com/VladimirReshetnikov/ProveIt/blob/21c0252ffe8611baa0fbf528fbb003d6c42ca9af/Analysis/GeneralizedRiemannHypothesis/Research/grh_transfers_lubin_tate_and_graph_towers/graphs.tex}{ProveIt formula}
gives $v_p(\kappa_n)=A(p^n-1)/(p-1)-2n$ for the spanning-tree
count of the corresponding Cayley tower.
Distinct nonzero signed column classes ensure simplicity at every layer.

\begin{theorem}[Three-step distribution]
For $m=3$ put $E=A-2(p+1)$ and $t=p-7$. The counts of row planes are:
\begin{center}
\begin{tabular}{crr}
\toprule
$E$ & $p\equiv1\bmod6$ & $p\equiv5\bmod6$\\
\midrule
0 & $p(p-1)/2$ & $p(p-1)/2$\\
2 & $t$ & $p+1$\\
4 & $8+\binom t2$ & $p(p+1)/2$\\
6 & $8t$ & $0$\\
8 & $28$ & $0$\\
\bottomrule
\end{tabular}
\end{center}
For the simple-tower counts remove $3$ planes from $E=4$ if
$-1$ is a square and from $E=0$ otherwise; also remove $6$ from
$E=4$ if $-2$ is a square and from $E=0$ otherwise.
\end{theorem}
\begin{proof}
Every $\PP(W)$ is a line in $\PP^2$ and cannot be contained in the
smooth conic $Q:\sum x_i^2=0$. The excess is two per conic point
and an extra two per point with $\sum x_i^4=0$. A nondegenerate
ternary form over a finite field is isotropic (compare the two
sets of square values in a diagonal equation); projection from
a rational conic point gives $p+1$ points. Thus there are
$p+1$ tangents, $\binom{p+1}{2}$ secants, and $p(p-1)/2$ external
lines. Newton's identities show that the special points have
squared coordinates the roots of $T^3-c$, $c\ne0$. They exist
exactly when $p\equiv1\bmod6$, in which case there are eight:
normalize the first coordinate to $1$, then order the other
two squared coordinates and choose their signs independently.
Counting tangents and unordered secant pairs by the number of
special points proves the table.

Nonsimple row planes have a normal proportional to a coordinate
vector or to $e_i\pm e_j$. There are three of the first type,
six of the second. Their lines impose $x_i=0$ or $x_i=\pm x_j$,
so contain no special points. Substituting into the conic gives
a secant or an external line according as $-1$ or $-2$ is a
square or a nonsquare. The simple-tower maximum thus agrees
with the unrestricted one for both congruence classes.
\end{proof}

For $p\equiv1\bmod6$, fix $z$ of order three and a special point
$a=(1,z,z^2)$. The signed permutation group is transitive on the
eight special points. The stabilizer of $a$ contains cyclic
permutation of the coordinates, acting projectively. On the other
seven points it fixes $(1,z^2,z)$ and cycles the three points of
each of the two remaining types: a single sign change of $a$, or a
single sign change of $(1,z^2,z)$. Taking cross products gives normals
\[
 (0,-z,1),\qquad (1,1,1),\qquad (z^2-z,1,-1)
\]
for the respective types (the middle case is the fixed point).
These three types cannot become equivalent even on unordered pairs:
the first normal has a zero coordinate, the second has three equal
squares, and the third has squares $-3,1,1$. The 28 maximizing planes
therefore form precisely three orbits, of sizes $12,4,12$.

\section{Four steps with vanishing quadratic moment}
Take $m=4$ and restrict to $M_1\equiv0$, i.e., maximal isotropic
planes $W$ for the dot product. The quadric
$Q:\sum_{i=1}^4 x_i^2=0$ is split: take $u^2+v^2=-1$ by
intersecting two sets of $(p+1)/2$ square values, then use the
orthogonal isotropic rows $(1,0,u,v),(0,1,-v,u)$. Its two rulings
therefore give $2(p+1)$ rational planes, and exactly two such
planes pass through each rational point of $Q$.
Put $C:\sum_i x_i^2=\sum_i x_i^4=0$ in $\PP^3$, and let
$c_p=\#C(\F_p)$.

\begin{theorem}[Four-step distribution and orbit]
For $m=4$ with $M_1=0$, set $E=A-4(p+1)$ and
$k=12\,{\bf 1}_{p\equiv1\bmod8}$. The numbers of row planes
at excess $0,8,16$ are respectively
\[
    2(p+1)-k-c_p/2,\qquad c_p/2,\qquad k.
\]
All these planes give simple towers. If $p\equiv1\bmod8$,
the twelve maximizing planes form one orbit under signed
permutations, also giving one orbit of maximizing residue
lists under change of basis, signs and permutations.
\end{theorem}

\begin{proof}
The quartic restricts nontrivially to every $W$ by the first-moment
lemma in the ProveIt report. Depth four requires that the squared
coordinates be a coset of the fourth roots of unity inside the
nonzero squares, by Newton's identities. Write $D$ for the set
of depth-four points. It is empty if $p\not\equiv1\bmod8$.
Otherwise, with $g$ of order eight, all $D$ points are signed
permutations of $a=(1,g,g^2,g^3)$ up to projective scaling.
Normalizing the first coordinate to $1$ gives
$3!\,2^3=48$ such points.

There are exactly two isotropic planes through $[a]$, obtained
by adjoining respectively $b=(-g,1,-g^3,g^2)$ and
$d=(-g,1,g^3,-g^2)$. They are distinct since their dot product is
$2(1+g^2)\ne0$. On the first plane the coordinate linear forms
are $U,V,g^2U,g^2V$ with $U=X-gY$, $V=gX+Y$. The quartic is
$2(U^4+V^4)$, with four distinct rational roots. The sextic
vanishes identically, so all four roots are in $D$.
On the second plane the quartic is $24g^2X^2Y^2$.
Both axis points are in $D$, and there are no other roots.
Consequently the planes that meet $D$ have either four $D$ points
and excess $16$, or two $D$ points and excess $8$.
Counting incidences with the 48 points of $D$ gives twelve of
the first type and twenty-four of the second. Transitivity on
$D$ and uniqueness of the four-point plane through each such
point prove the orbit assertion.

For completeness, on a plane avoiding $D$ the quartic has
either zero or four rational zeros. Indeed, make two independent
columns the coordinate basis and write the column matrix as
$[I\mid B]$. Since $BB^{\mathsf T}=-I$, a sign choice for the last
column puts $B$ in the form
$\left(\begin{smallmatrix}u&v\\v&-u\end{smallmatrix}\right)$.
With $J=\left(\begin{smallmatrix}0&-1\\1&0\end{smallmatrix}\right)$,
the matrices $B,J$ permute the signed columns and generate a
Klein four group projectively: $J^2=B^2=(JB)^2=-I$ and $JB=-BJ$.
On covectors the transposes preserve every moment. If a quartic
zero is fixed by a nonidentity group element, its covector
eigenvalue $\lambda$ has $\lambda^2=-1$ and hence
$\lambda^6=-1$. Its sextic vanishes, which would put the point
in $D$. Thus every quartic root on a plane avoiding $D$ belongs
to an orbit of four. Each such orbit gives excess $8$.

Every point of $C\setminus D$ occurs on two planes avoiding $D$.
It follows that $(c_p-\#D)/2$ such planes contribute excess $8$,
in addition to the $\#D/2$ planes having two depth-four points.
This gives the claimed counts. A nonsimple plane would lie in
the orthogonal complement of a coordinate vector or
$e_i\pm e_j$. Both complements are nondegenerate of dimension
three and cannot contain a totally isotropic plane.
\end{proof}

When $p\equiv1\bmod8$, lists with $M_1\ne0$ have at most two
exceptional directions, so $A\le2(p+1)+12$; they cannot be
maximizers. Outside that congruence the theorem gives the number
of maximizing row planes as $c_p/2$ if $c_p>0$, and $2(p+1)$
otherwise. The already known four-step maximum and the seven
primes with $c_p=0$ are in the linked ProveIt report. This note
adds the distribution and orbit statements, without an RH or
GRH assumption.
\end{document}
